\documentclass[12pt,a4paper]{article}

\usepackage[utf8]{inputenc}
\usepackage[T1]{fontenc}
\usepackage[brazilian]{babel}
\usepackage{newtxtext}
\usepackage{newtxmath}
\usepackage{geometry}
\usepackage{setspace}
\usepackage{array}
\usepackage{booktabs}
\usepackage{longtable}
\usepackage{tabularx}
\usepackage{xcolor}
\usepackage{enumitem}
\usepackage{hyperref}
\usepackage{titlesec}
\usepackage{fancyhdr}

\geometry{left=3cm,right=2cm,top=3cm,bottom=2cm}
\setstretch{1.5}
\setlength{\parindent}{1.25cm}
\setlength{\parskip}{0pt}
\pagestyle{fancy}
\fancyhf{}
\fancyhead[R]{\thepage}
\renewcommand{\headrulewidth}{0pt}
\setlength{\headheight}{14.5pt}
\hypersetup{
    colorlinks=true,
    linkcolor=black,
    urlcolor=black,
    pdftitle={Projeto Aplicado II - Etapa 1},
    pdfauthor={Grupo de Trabalho},
    pdfsubject={Apresentacao, objetivos, metas e milestones do projeto}
}

\definecolor{maincolor}{RGB}{0, 0, 0}

\titleformat{\section}{\large\bfseries\color{maincolor}}{\thesection.}{0.5em}{}
\titleformat{\subsection}{\normalsize\bfseries\color{maincolor}}{\thesubsection.}{0.5em}{}

\newcolumntype{Y}{>{\raggedright\arraybackslash}X}

\begin{document}
\hypersetup{pageanchor=false}

\begin{titlepage}
    \centering
    {\large Universidade Presbiteriana Mackenzie \par}
    \vspace{0.4cm}
    {\large Banco de Dados – Projeto Aplicado II \par}
    \vspace{2.0cm}
    {\Large \textbf{Etapa 1} \par}
    \vspace{0.4cm}
    {\large \textbf{Apresentação, Objetivos, Metas e Milestones do Projeto} \par}
    \vspace{1.4cm}
    {\Large \textbf{Classificação e Monitoramento Textual de Boletins Epidemiológicos sobre Dengue no Brasil} \par}
    \vfill

    \begin{flushleft}
    \textbf{Grupo:} Ana Clara Silva de Souza; Cid Wallace Araujo de Oliveira; Eduardo Machado Silva; Frederico Ripamonte Borges \\
    \textbf{Disciplina:} Projeto Aplicado II \\
    \textbf{Professor:} Anderson Adaime de Borba \\
    \textbf{Semestre:} 2026.01
    \end{flushleft}

    \vfill
    {\large \today \par}
\end{titlepage}

\clearpage
\hypersetup{pageanchor=true}
\setcounter{page}{1}

\section*{Resumo Executivo}

Este documento apresenta a proposta inicial do projeto da disciplina \textit{Projeto Aplicado II}, cujo foco é a aplicação de técnicas de ciência de dados e análise textual sobre documentos públicos relacionados à dengue no Brasil. O projeto utilizará como principal fonte institucional o Ministério da Saúde, a partir de boletins epidemiológicos, informes técnicos, comunicados oficiais e demais documentos textuais disponibilizados em repositórios públicos governamentais.

A proposta parte de um problema prático: há grande volume de informação textual relevante sobre dengue, mas sua leitura, organização e interpretação manual são lentas, dispersas e pouco escaláveis. Nesse contexto, o projeto pretende estruturar uma base textual com metadados padronizados e preparar o caminho para o desenvolvimento de um produto analítico capaz de classificar e organizar automaticamente conteúdos temáticos relevantes ao monitoramento em saúde pública.

A presente entrega contempla os itens exigidos na Etapa 1: definição do grupo de trabalho, premissas do projeto, definição da organização escolhida, área de atuação, apresentação dos dados que serão utilizados, objetivos, metas, milestones e cronograma estimado de execução.

\clearpage
\section{Grupo de Trabalho}

O projeto será desenvolvido em grupo, conforme exigência da disciplina. A composição final deve conter de três a quatro integrantes.

\subsection{Integrantes}

\begin{itemize}[leftmargin=1.2cm]
    \item Ana Clara Silva de Souza
    \item Cid Wallace Araujo de Oliveira
    \item Eduardo Machado Silva
    \item Frederico Ripamonte Borges
\end{itemize}

\subsection{Divisão inicial de responsabilidades}

Para favorecer a execução colaborativa e a organização do projeto, propõe-se a seguinte divisão inicial de responsabilidades, sem prejuízo de atuação conjunta em todas as etapas:

\begin{itemize}[leftmargin=1.2cm]
    \item \textbf{Coordenação e documentação:} organização do cronograma, consolidação das entregas e redação dos relatórios técnicos;
    \item \textbf{Coleta e estruturação de dados:} levantamento das fontes oficiais, captura dos documentos e organização da base textual;
    \item \textbf{Tratamento e análise de dados:} preparação textual, modelagem e avaliação analítica nas etapas posteriores;
    \item \textbf{Visualização e apresentação:} elaboração de tabelas, gráficos, storytelling e apresentação final.
\end{itemize}

Essa divisão tem caráter operacional e poderá ser ajustada ao longo do semestre, de acordo com a evolução do projeto e a disponibilidade do grupo.

\clearpage
\section{Premissas do Projeto}

\subsection{Tema do Projeto}

O tema definido pelo grupo é a \textbf{análise textual de documentos epidemiológicos sobre dengue no Brasil}, com foco em organização, classificação temática e extração de informações úteis ao monitoramento em saúde pública.

\subsection{Organização Escolhida}

A organização escolhida como referência institucional do projeto é o \textbf{Ministério da Saúde do Brasil}, por meio de suas bases públicas e documentos oficiais relacionados ao acompanhamento da dengue no país.

A escolha dessa organização se justifica por quatro fatores principais. Em primeiro lugar, trata-se de uma fonte pública, oficial e relevante socialmente. Em segundo, o Ministério da Saúde disponibiliza documentos textuais com alto valor informacional, adequados à proposta da disciplina, que exige manipulação de texto ou imagem. Em terceiro, o tema apresenta forte aderência a aplicações reais de ciência de dados em saúde pública. Por fim, a utilização de dados governamentais aumenta a transparência metodológica e a reprodutibilidade do trabalho.

\subsection{Área de Atuação}

A área de atuação do projeto é definida como:

\begin{quote}
\textbf{Aplicação de técnicas de ciência de dados e análise textual em documentos epidemiológicos sobre dengue, com foco na organização, classificação e extração de informações úteis ao monitoramento em saúde pública.}
\end{quote}

Essa definição delimita com clareza o escopo do trabalho: o projeto não buscará apenas descrever documentos, mas estruturar uma base textual capaz de sustentar análises posteriores e a construção de um produto analítico.

\subsection{Problema de Interesse}

A dengue é uma arbovirose de grande relevância no Brasil, acompanhada por múltiplos documentos institucionais que incluem alertas epidemiológicos, orientações técnicas, informações clínicas, ações preventivas e dados consolidados. Apesar da abundância de material oficial, o acesso rápido às informações mais relevantes pode ser dificultado pela dispersão dos conteúdos, pela heterogeneidade dos formatos e pelo volume documental acumulado ao longo do tempo.

Diante disso, o problema central do projeto pode ser resumido da seguinte forma:

\begin{quote}
\textit{Como organizar e analisar documentos públicos sobre dengue de maneira padronizada, de modo a facilitar a identificação de temas, padrões e informações relevantes para monitoramento em saúde pública?}
\end{quote}

\subsection{Proposta de Solução}

A proposta do grupo é construir, inicialmente, uma base textual estruturada a partir de documentos oficiais sobre dengue. Em etapas posteriores, essa base servirá de insumo para o desenvolvimento de um produto analítico de classificação textual, capaz de categorizar automaticamente conteúdos segundo temas de interesse, tais como cenário epidemiológico, prevenção, orientações clínicas, campanhas ou recomendações operacionais.

Assim, já na Etapa 1, o projeto é concebido de forma integrada e progressiva: a definição da fonte, dos dados, dos objetivos e do cronograma foi feita considerando a viabilidade das entregas futuras.

\clearpage
\section{Apresentação dos Dados que Serão Utilizados}

\subsection{Natureza dos Dados}

A base principal do projeto será composta por \textbf{dados textuais}, extraídos de documentos públicos oficiais relacionados à dengue. Entre os tipos de documentos que poderão compor o corpus, destacam-se:

\begin{itemize}[leftmargin=1.2cm]
    \item boletins epidemiológicos;
    \item informes técnicos;
    \item notas e comunicados oficiais;
    \item páginas institucionais informativas;
    \item documentos textuais complementares publicados em portais governamentais de saúde.
\end{itemize}

O corpus será constituído prioritariamente por materiais vinculados ao Ministério da Saúde e, quando pertinente, por conteúdos oficiais correlatos que mantenham aderência ao escopo do projeto.

\subsection{Unidade de Análise}

Para esta etapa inicial, a unidade documental primária será o \textbf{documento oficial} como um todo. Contudo, tendo em vista as etapas posteriores do projeto, prevê-se também a segmentação dos documentos em \textbf{trechos, seções ou parágrafos}, quando isso for necessário para treinamento e avaliação de métodos analíticos.

Essa decisão é estratégica, pois amplia a flexibilidade do projeto: permite iniciar com uma base documental consistente e, ao mesmo tempo, preparar o material para tarefas futuras de classificação textual.

\subsection{Estrutura Proposta dos Metadados}

Com o objetivo de garantir organização, rastreabilidade e reprodutibilidade, cada documento deverá ser registrado com um conjunto padronizado de metadados.

\begin{center}
\renewcommand{\arraystretch}{1.25}
\begin{tabularx}{\textwidth}{>{\bfseries}p{4.5cm} Y}
\toprule
Campo & Descrição \\
\midrule
ID do documento & Identificador único para controle interno da base \\
Título & Nome oficial do documento \\
Data de publicação & Data em que o documento foi disponibilizado \\
Órgão responsável & Instituição responsável pela publicação \\
Tipo de documento & Boletim, informe técnico, comunicado, página institucional etc. \\
Tema principal & Assunto predominante do documento \\
Período de referência & Intervalo temporal abordado no conteúdo, quando houver \\
Escopo geográfico & Brasil, região, estado ou outro recorte indicado no documento \\
Fonte / URL & Endereço eletrônico da publicação original \\
Texto bruto & Conteúdo extraído originalmente \\
Texto tratado & Conteúdo após limpeza e padronização \\
Palavras-chave & Termos relevantes identificados no documento \\
Observações & Informações adicionais úteis para análise \\
\bottomrule
\end{tabularx}
\end{center}

\subsection{Exemplo de Uso dos Metadados}

Os metadados terão papel central na organização e no uso analítico da base. Eles permitirão:

\begin{itemize}[leftmargin=1.2cm]
    \item localizar rapidamente documentos por tema, período ou fonte;
    \item padronizar a estrutura do corpus textual;
    \item facilitar a preparação dos dados para as etapas de tratamento e modelagem;
    \item garantir rastreabilidade entre o conteúdo analisado e sua fonte original;
    \item apoiar visualizações e análises descritivas sobre a composição da base.
\end{itemize}

\subsection{Justificativa para a Escolha dos Dados}

A escolha dos dados textuais oficiais é adequada à proposta do projeto por três razões principais. Primeiro, os documentos apresentam relevância pública e aderência temática. Segundo, o material textual permite aplicar técnicas de preparação de dados e aprendizado de máquina vistas ao longo do semestre. Terceiro, trata-se de uma base com potencial real de transformação em produto analítico, o que fortalece a coerência entre as quatro etapas do projeto.

\clearpage
\section{Objetivos do Projeto}

\subsection{Objetivo Geral}

Desenvolver uma base textual estruturada e um protótipo de análise/classificação de documentos epidemiológicos sobre dengue no Brasil, utilizando dados públicos oficiais para organizar, extrair e interpretar informações relevantes ao monitoramento em saúde pública.

\subsection{Objetivos Específicos}

\begin{enumerate}[leftmargin=1.2cm]
    \item Levantar e selecionar documentos públicos oficiais relacionados à dengue;
    \item Organizar os documentos em uma base padronizada com metadados;
    \item Realizar a limpeza e a preparação textual da base;
    \item Definir categorias temáticas adequadas ao corpus;
    \item Aplicar métodos analíticos de classificação textual em etapas posteriores;
    \item Avaliar o desempenho analítico por métricas de acurácia;
    \item Consolidar os resultados em visualizações, relatório técnico e apresentação final.
\end{enumerate}

\clearpage
\section{Metas do Projeto}

As metas a seguir foram definidas de forma realista e progressiva, buscando alinhar qualidade técnica, viabilidade de execução e aderência aos requisitos da disciplina.

\begin{enumerate}[leftmargin=1.2cm]
    \item Estruturar um repositório colaborativo no GitHub para versionamento do projeto;
    \item Coletar e catalogar uma base inicial de documentos oficiais sobre dengue;
    \item Definir um padrão mínimo de metadados para todo o corpus;
    \item Preparar a base textual para análise exploratória e treinamento;
    \item Comparar, ao menos, dois métodos analíticos aplicáveis ao problema proposto;
    \item Apresentar um produto analítico preliminar na Etapa 3;
    \item Entregar, ao final, relatório técnico, storytelling, repositório completo e apresentação gravada.
\end{enumerate}

\subsection{Metas Mensuráveis}

Para aumentar a objetividade do acompanhamento, o grupo estabelece as seguintes metas mensuráveis iniciais:

\begin{itemize}[leftmargin=1.2cm]
    \item catalogar ao menos 30 documentos textuais oficiais na base inicial;
    \item adotar no mínimo 10 campos padronizados de metadados;
    \item estruturar pelo menos 1 base consolidada em formato tabular para análise;
    \item testar pelo menos 2 abordagens analíticas na etapa de modelagem;
    \item produzir pelo menos 1 visualização analítica e 1 protótipo demonstrável antes da entrega final.
\end{itemize}

Esses números poderão ser ajustados conforme disponibilidade e qualidade das fontes, sem alteração da lógica central do projeto.

\clearpage
\section{Milestones do Projeto}

Os principais \textit{milestones} definidos pelo grupo são os seguintes:

\begin{center}
\renewcommand{\arraystretch}{1.25}
\begin{tabularx}{\textwidth}{>{\bfseries}p{3.0cm} Y}
\toprule
Milestone & Descrição \\
\midrule
M1 & Tema, organização escolhida, objetivos, metas e cronograma definidos \\
M2 & Base inicial de documentos coletada e repositório colaborativo ativo \\
M3 & Metadados estruturados e corpus textual preparado para exploração \\
M4 & Método analítico escolhido e base pronta para treinamento \\
M5 & Resultados preliminares consolidados, com métricas e produto inicial \\
M6 & Relatório final, storytelling, GitHub completo e vídeo de apresentação entregues \\
\bottomrule
\end{tabularx}
\end{center}

\clearpage
\section{Cronograma de Atividades}

O cronograma abaixo foi construído considerando as quatro etapas formais da disciplina e as datas previstas de entrega.

{\small
\setstretch{1.15}
\begin{center}
\renewcommand{\arraystretch}{1.08}
\begin{tabularx}{\textwidth}{>{\bfseries}p{1.6cm} >{\raggedright\arraybackslash}p{2.5cm} Y Y}
\toprule
Etapa & Período estimado & Atividades principais & Entrega prevista \\
\midrule
Etapa 1 & 16/03/2026 & Definição do grupo, organização, problema, dados, objetivos, metas, milestones e cronograma & Relatório técnico inicial \\
\addlinespace
Etapa 2 & 30/03/2026 & Implantação do GitHub, definição de bibliotecas, análise exploratória, tratamento dos dados e definição do método analítico e da acurácia & Relatório técnico da definição do produto analítico \\
\addlinespace
Etapa 3 & 27/04/2026 & Aplicação do método analítico, consolidação dos resultados, cálculo das métricas, apresentação do produto e esboço do storytelling & Relatório técnico com resultados preliminares e storytelling inicial \\
\addlinespace
Etapa 4 & 25/05/2026 & Consolidação final, refinamento do storytelling, organização do GitHub, preparação da apresentação e gravação do vídeo & Relatório técnico final, apresentação, GitHub e vídeo \\
\bottomrule
\end{tabularx}
\end{center}

\subsection{Plano Operacional Resumido}

Além da divisão por etapas, o grupo adotará um fluxo contínuo de trabalho:

\begin{itemize}[leftmargin=1.2cm,itemsep=0.2em]
    \item reuniões periódicas para acompanhamento das atividades;
    \item versionamento de código e documentos no GitHub;
    \item validação interna de cada entrega antes da submissão;
    \item registro de decisões metodológicas ao longo do semestre.
\end{itemize}
}

\clearpage
\section{Link do GitHub do Projeto}

O repositório do projeto será utilizado como ambiente central de colaboração, versionamento e organização dos arquivos produzidos pelo grupo.

\textbf{Link do GitHub:}\\
\url{https://github.com/fredericorbgs/projeto-aplicado-ii-dengue}

\subsection{Estrutura inicial sugerida do repositório}

\begin{itemize}[leftmargin=1.2cm]
    \item \texttt{README.md}
    \item \texttt{data/}
    \item \texttt{notebooks/}
    \item \texttt{src/}
    \item \texttt{docs/}
    \item \texttt{reports/}
\end{itemize}

Essa organização facilitará tanto a execução colaborativa quanto a rastreabilidade das entregas técnicas ao longo do semestre.

\clearpage
\section{Considerações Finais}

A proposta apresentada nesta Etapa 1 foi estruturada para atender de forma completa aos requisitos da disciplina e, ao mesmo tempo, estabelecer uma base sólida para as próximas entregas. A escolha do Ministério da Saúde como fonte institucional, aliada ao uso de documentos textuais oficiais sobre dengue, fornece ao projeto relevância pública, coerência metodológica e viabilidade técnica.

Além de cumprir os requisitos imediatos da etapa inicial, o projeto já nasce com uma direção analítica clara: transformar um conjunto disperso de documentos em uma base padronizada, explorável e apta a sustentar técnicas de classificação textual e geração de insights. Dessa forma, o grupo pretende conduzir o trabalho com consistência entre planejamento, execução técnica e comunicação dos resultados.

\vspace{1cm}

% --------------------------------------------------------------------
% Somatoria: inclui a Etapa 2 (A2) ao final do documento da Etapa 1.
\clearpage
\hypersetup{pageanchor=false}

\begin{titlepage}
    \centering
    {\large Universidade Presbiteriana Mackenzie \par}
    \vspace{0.4cm}
    {\large Banco de Dados – Projeto Aplicado II \par}
    \vspace{2.0cm}
    {\Large \textbf{Etapa 2} \par}
    \vspace{0.4cm}
    {\large \textbf{Definição do Método Analítico} \par}
    \vspace{1.4cm}
    {\Large \textbf{Classificação e Monitoramento Textual de Boletins Epidemiológicos sobre Dengue no Brasil} \par}
    \vfill

    \begin{flushleft}
    \textbf{Grupo:} Ana Clara Silva de Souza; Cid Wallace Araujo de Oliveira; Eduardo Machado Silva; Frederico Ripamonte Borges \\
    \textbf{Disciplina:} Projeto Aplicado II \\
    \textbf{Professor:} Anderson Adaime de Borba \\
    \textbf{Semestre:} 2026.01
    \end{flushleft}

    \vfill
    {\large \today \par}
\end{titlepage}

\clearpage
\hypersetup{pageanchor=true}
\setcounter{page}{1}

\section*{Resumo Executivo}

Esta segunda etapa do projeto apresenta a definição do método analítico que será utilizado para a classificação textual supervisionada de boletins epidemiológicos relacionados a arboviroses no Brasil. Em coerência com a proposta da Etapa 1, a pesquisa baseia-se no uso de técnicas de processamento de linguagem natural para estruturar uma base analítica e treinar modelos com representações numéricas dos textos.

Nesta entrega, o grupo define a linguagem de programação empregada (Python), detalha como a base foi analisada exploratoriamente (EDA), descreve o tratamento previsto para extração, limpeza, tokenização e segmentação (chunking) para preparação do conjunto de treino e avaliação. Adicionalmente, são apresentadas as bases teóricas dos métodos utilizados (TF-IDF, Naive Bayes Multinomial e Regressão Logística multiclasse) e a forma de cálculo da métrica principal de avaliação (acurácia), incluindo complementos para leitura em cenários multiclasse.

% --------------------------------------------------------------------
\clearpage
\section{Definição da Linguagem de Programação}

\subsection{Linguagem Escolhida}

A linguagem de programação escolhida para o desenvolvimento do projeto é o \textbf{Python}.

\subsection{Justificativa da Escolha}

A decisão por Python se justifica por sua adequação ao escopo do projeto, que envolve manipulação de dados e processamento de texto. Em termos práticos, Python permite:

\begin{enumerate}[leftmargin=1.2cm]
    \item Construir um pipeline reprodutível com leitura/escrita de artefatos (por exemplo: PDFs, textos extraídos, chunks e CSVs consolidados).
    \item Aplicar rotinas de extração textual a partir de PDFs, quando necessário, por meio de bibliotecas específicas.
    \item Representar e transformar textos em vetores numéricos por meio de \textit{TF-IDF} e treinar classificadores supervisionados com \texttt{scikit-learn}.
    \item Avaliar modelos com métricas adequadas para classificação multiclasse, como acurácia e métricas complementares (precision, recall, F1 e matriz de confusão).
    \item Organizar a implementação em módulos claros (extração, pré-processamento, modelagem e avaliação), facilitando a colaboração do grupo e a rastreabilidade das decisões.
\end{enumerate}

\subsection{Bibliotecas Previstas}

As bibliotecas usadas/previstas no desenvolvimento do pipeline incluem:

\begin{table}[h]
\centering
\renewcommand{\arraystretch}{1.25}
\begin{tabularx}{\textwidth}{>{\bfseries}p{3.6cm} Y}
\toprule
Biblioteca & Finalidade no projeto \\
\midrule
\texttt{pandas} & Organização tabular da base, leitura de CSVs e estruturação de metadados \\
\texttt{numpy} & Operações numéricas auxiliares \\
\texttt{matplotlib}/\texttt{seaborn} & Visualizações para EDA \\
\texttt{scikit-learn} & TF-IDF, modelos (Naive Bayes, Regressão Logística) e avaliação \\
\texttt{requests} & Coleta automatizada (etapa de coleta, quando aplicável) \\
\texttt{pdfplumber} & Extração de texto de PDFs \\
\bottomrule
\end{tabularx}
\caption{Bibliotecas utilizadas/previstas para o projeto}
\label{tab:bibliotecas}
\end{table}

% --------------------------------------------------------------------
\clearpage
\section{Análise Exploratória da Base de Dados}

\subsection{Base Escolhida}

A base principal do projeto é composta por boletins epidemiológicos em \textbf{PDF} relacionados a arboviroses urbanas (incluindo Dengue, Zika e Chikungunya). Para esta etapa, os boletins utilizados foram extraídos do seguinte endereço oficial:

\vspace{0.2cm}
\noindent
\textbf{Fonte oficial dos boletins (COVISA/SMS-SP):} \\
\url{https://prefeitura.sp.gov.br/web/saude/w/vigilancia_em_saude/boletim_covisa/arboviroses}

\vspace{0.2cm}
\noindent
Observação do próprio site: recomenda-se utilizar o \textit{último boletim publicado}, e boletins anteriores devem ser entendidos como uma fotografia do momento em que foram publicados (não é recomendada comparação direta entre anos diferentes com o boletim vigente).

\subsection{Objetivo da Análise Exploratória}

A análise exploratória tem como objetivo compreender a estrutura do corpus textual, identificar padrões relevantes para modelagem e orientar a preparação dos dados (limpeza, tokenização e segmentação). Além da descrição, a EDA serve para verificar:

\begin{enumerate}[leftmargin=1.2cm]
    \item Quantidade de documentos disponíveis e distribuição de tamanho.
    \item Distribuição de vocabulário e termos predominantes.
    \item Necessidade de segmentação do texto (evitar mistura temática em uma única amostra).
    \item Presença de ruídos típicos da extração de PDFs e como mitigá-los.
    \item Parâmetros de chunking adequados para viabilizar treino e avaliação supervisionada.
\end{enumerate}

\subsection{Unidade de Análise}

A unidade de análise primária é o documento (boletim completo). Entretanto, para aumentar a capacidade do modelo em capturar temas locais e reduzir mistura temática, a etapa prevê a segmentação do texto em \textbf{chunks} (trechos) a partir de janelas de tokens.

\subsection{Resultados Reais da EDA (base processada)}

Com os boletins disponíveis, foram observados os seguintes resultados após extração, limpeza e preparação do texto (conforme artefatos em \texttt{data/processed/}):

\begin{itemize}[leftmargin=1.2cm]
    \item \textbf{Quantidade de PDFs:} 11
    \item \textbf{Chunks gerados:} 60
    \item \textbf{Parâmetros de chunking:} \texttt{chunk\_size=500}, \texttt{stride=200}, \texttt{min\_fill\_ratio=0.6}
    \item \textbf{Tokens por chunk:} média 487.62, mínimo 356, máximo 500
    \item \textbf{Vocabulário (termos únicos após filtragem):} 440
\end{itemize}

\subsubsection{Termos predominantes}

Os termos mais frequentes do corpus, após filtragem (remoção de números puros e tokens muito curtos), reforçam que o conteúdo é dominado por vocabulário epidemiológico e pela estrutura recorrente dos boletins:

\begin{table}[h]
\centering
\renewcommand{\arraystretch}{1.15}
\begin{tabularx}{\textwidth}{>{\bfseries}l Y}
\toprule
Termo & Ocorrências \\
\midrule
casos & 799 \\
sao & 669 \\
ate & 600 \\
confirmados & 569 \\
sinan & 550 \\
municipio & 509 \\
inicio & 507 \\
quadro & 502 \\
online & 494 \\
ano & 485 \\
sintomas & 485 \\
entre & 476 \\
vila & 461 \\
paulo & 436 \\
segundo & 432 \\
\bottomrule
\end{tabularx}
\caption{Top termos mais frequentes (EDA)}
\label{tab:toptermos}
\end{table}

\subsubsection{Implicações para a modelagem}

Com base na EDA, o grupo identifica as seguintes implicações práticas:

\begin{enumerate}[leftmargin=1.2cm]
    \item A recorrência de termos sugere que a representação por \textbf{TF-IDF} é adequada para capturar termos discriminantes entre categorias.
    \item A presença de conteúdo relacionado a mais de uma arbovirose no mesmo boletim aumenta a chance de mistura temática; por isso, a segmentação em chunks favorece a qualidade das amostras.
    \item A similaridade entre boletins (formatos e seções recorrentes) motiva uma estratégia de rotulagem por trechos (chunks), reduzindo mistura de temas em uma única amostra.
\end{enumerate}

\subsubsection{Ruído da extração e limpeza aplicada}

Como PDFs podem introduzir artefatos na extração textual, foram aplicadas rotinas de limpeza antes do chunking:

\begin{itemize}[leftmargin=1.2cm]
    \item Uniao de hifens de quebra de linha (ex.: \texttt{transmissao-continua} $\rightarrow$ \texttt{transmissaocontinua}).
    \item Remoção de linhas que se tornam apenas números (ex.: numeração de páginas/rodapés).
    \item Normalização de \textit{whitespace} (remoção de múltiplos espaços e excesso de quebras de linha).
    \item Normalização de acentos para ASCII (portabilidade dos artefatos no repositório).
    \item Tokenização por \textit{regex} e filtragem (remoção de números puros e tokens com menos de 3 caracteres).
\end{itemize}

% --------------------------------------------------------------------
\clearpage
\section{Tratamento da Base de Dados}

\subsection{Visão Geral do Pipeline}

O tratamento da base de dados compreende um pipeline composto por:

\begin{enumerate}[leftmargin=1.2cm]
    \item Coleta/importação dos PDFs para a estrutura do repositório.
    \item Extração textual a partir dos PDFs.
    \item Limpeza e normalização do texto.
    \item Tokenização e filtragem para compatibilidade com TF-IDF.
    \item Segmentação (chunking) em janelas de tokens.
    \item Criação de um dataset consolidado (CSV) para modelagem.
    \item Rotulagem supervisionada (quando as categorias forem definidas pelo grupo).
    \item Vetorização e treino dos modelos.
    \item Avaliação das predições no conjunto de teste.
\end{enumerate}

\subsection{Preparação (extração, limpeza, tokenização e chunking)}

O processo foi implementado no módulo \texttt{src/preprocessing.py} e gera os seguintes artefatos:

\begin{itemize}[leftmargin=1.2cm]
    \item \texttt{data/processed/extracted\_text/}: texto limpo por documento (\texttt{.md}).
    \item \texttt{data/processed/chunks/}: texto por chunk (\texttt{.md}).
    \item \texttt{data/processed/chunks\_dataset.csv}: dataset consolidado para modelagem.
    \item \texttt{data/processed/eda\_summary.*}: resumo estatístico da EDA (tokens, vocabulário e top termos).
\end{itemize}

\subsubsection{Chunking e identificação de amostras}

Cada chunk recebe um identificador no formato:

\[
 \text{chunk\_id} = \text{doc\_id}\_\text{chunk\_index}
 \]

Com:
\begin{itemize}[leftmargin=1.2cm]
    \item \texttt{doc\_id}: derivado do nome do arquivo PDF (sem extensão).
    \item \texttt{chunk\_index}: índice sequencial do chunk dentro do documento.
\end{itemize}

\subsubsection{Dataset final para modelagem}

O arquivo \texttt{data/processed/chunks\_dataset.csv} contém (no mínimo) as colunas:

\begin{itemize}[leftmargin=1.2cm]
    \item \texttt{chunk\_id}
    \item \texttt{doc\_id}
    \item \texttt{chunk\_index}
    \item \texttt{token\_count}
    \item \texttt{chunk\_text}
    \item \texttt{text\_md\_path}
\end{itemize}

\subsection{Rotulagem supervisionada (labels)}

Para que o treino supervisionado ocorra, será necessário criar:

\begin{itemize}[leftmargin=1.2cm]
    \item \texttt{data/processed/labels.csv} com colunas:
    \begin{itemize}
        \item \texttt{chunk\_id}
        \item \texttt{label}
    \end{itemize}
\end{itemize}

\subsection{Vetorização e Treinamento}

Com o dataset rotulado, os textos serão representados numericamente por TF-IDF e treinados em dois classificadores:

\begin{enumerate}[leftmargin=1.2cm]
    \item \textbf{Baseline (referência inicial)}: \texttt{MultinomialNB} (\textit{Naive Bayes Multinomial}).
    \item \textbf{Modelo principal}: \texttt{LogisticRegression} em modo multiclasse (\textit{multinomial}) para classificação discriminativa.
\end{enumerate}

\subsubsection{Parâmetros de TF-IDF}

O \texttt{TfidfVectorizer} utiliza:

\begin{itemize}[leftmargin=1.2cm]
    \item \texttt{ngram\_range=(1,2)} (unigramas e bigramas)
    \item \texttt{min\_df=2} e \texttt{max\_df=0.95}
    \item \texttt{max\_features=20000}
    \item \texttt{stop\_words=None} (mantém termos do domínio)
\end{itemize}

\subsubsection{Divisão de dados (treino/validação/teste)}

Quando labels estiverem disponíveis, a divisão será estratificada por classe:

\begin{itemize}[leftmargin=1.2cm]
    \item 70\% treino
    \item 15\% validação
    \item 15\% teste final
\end{itemize}

% --------------------------------------------------------------------
\clearpage
\section{Definição e Descrição das Bases Teóricas dos Métodos}

\subsection{Aprendizado Supervisionado para Classificação Textual}

O problema será formulado como uma tarefa de \textbf{classificação supervisionada}: os textos vetorizados são associados a rótulos (categorias temáticas previamente atribuídas). O objetivo é aprender padrões que permitam generalizar a classificação para novos textos.

\subsection{Base Teórica da Representação TF-IDF}

A técnica TF-IDF atribui peso maior a termos que aparecem com frequência no documento, mas que não são excessivamente comuns no corpus. Formalmente:

\[
\operatorname{TF\mbox{-}IDF}(t,d) = \operatorname{tf}(t,d)\cdot \log\left(\frac{N}{\operatorname{df}(t)}\right)
\]

em que:
\begin{itemize}[leftmargin=1.2cm]
    \item $N$ é o número total de documentos no corpus;
    \item $\operatorname{df}(t)$ é o número de documentos em que o termo $t$ aparece;
    \item $\operatorname{tf}(t,d)$ é a frequência do termo $t$ no documento $d$.
\end{itemize}

\subsection{Modelo 1: Naive Bayes Multinomial}

O \textit{Naive Bayes Multinomial} é um classificador probabilístico baseado na probabilidade de uma classe dado o conjunto de palavras observadas no texto. Sua decisão é expressa como:

\[
\hat{c} = \arg\max_{c \in C} P(c)\prod_{i=1}^{n}P(w_i \mid c)
\]

\subsection{Modelo 2: Regressão Logística Multiclasse}

A \textit{Regressão Logística Multiclasse} estima probabilidades por classe. Com função softmax:

\[
P(y = k \mid x) = \frac{e^{\beta_k^\top x}}{\sum_{j=1}^{K} e^{\beta_j^\top x}}
\]

\subsection{Estratégia Analítica Adotada}

O grupo utiliza o \textit{Naive Bayes Multinomial} como modelo de referência e a \textit{Regressão Logística Multiclasse} como modelo principal para comparação.

% --------------------------------------------------------------------
\clearpage
\section{Definição e Descrição de Como Será Calculada a Acurácia}

\subsection{Métrica Principal}

A métrica principal de avaliação será a \textbf{acurácia}, definida como a proporção de previsões corretas em relação ao total de previsões:

\[
\text{Acurácia} = \frac{\text{número de classificações corretas}}{\text{número total de classificações}}
\]

Essa métrica será calculada no conjunto de teste (com rótulos reais) após o treino e a seleção usando apenas a estratégia de validação, evitando vazamento de informação.

\subsection{Métricas Complementares}

Além da acurácia, o grupo calculará métricas complementares para classificação multiclasse:
\begin{itemize}[leftmargin=1.2cm]
    \item \textbf{Precision} por classe e média macro;
    \item \textbf{Recall} por classe e média macro;
    \item \textbf{F1-score} por classe e média macro;
    \item \textbf{Matriz de confusão} para análise de erros entre categorias.
\end{itemize}

% --------------------------------------------------------------------
\clearpage
\section{Considerações Finais}

Com base na EDA realizada sobre boletins epidemiológicos em PDF (11 documentos, segmentados em 60 chunks), o grupo definiu uma abordagem metodológica coerente para classificação textual supervisionada. O pipeline de tratamento (extração, limpeza, tokenização, chunking e construção do dataset) fornece uma base reprodutível para o treino de modelos via TF-IDF.

\vspace{0.3cm}
\noindent
\textbf{Fonte oficial dos boletins:} \\
\url{https://prefeitura.sp.gov.br/web/saude/w/vigilancia_em_saude/boletim_covisa/arboviroses}



\end{document}